An architecture can have enough parameters to represent a relational map and still lack the
information contract required by a particular interpretation.  A route selected after reading both
endpoint residues is a valid adaptive computation.  Holding that same route fixed during an
endpoint intervention requires a context that can be reconstructed without those endpoint values.

This chapter develops the resulting information boundary.  Endpoint-removable state gives one
constructive mechanism.  Strict targetwise support states the measurable object.  A spanning-support
theorem limits globally invariant adaptive graphs, and a Bayes floor quantifies the prediction error
that remains when endpoint evidence is withheld.  These results continue the endpoint-exclusion
condition introduced in Chapter~\ref{chap:book-static-dynamic} and separate semantic realizability
from the capacity and execution questions of Chapter~\ref{chap:book-architecture}.
